\documentclass[twocolumn]{iopjournal}

\usepackage{amsmath,amssymb,amsfonts}

\fancyhfoffset[L]{0mm}
\setlength{\marginparwidth}{0pt}
\setlength{\marginparsep}{0pt}

% ==========================================
% ZENODO PREPRINT HEADERS (ACTIVE)
% ==========================================
\renewcommand{\articletype}[1]{%
  {\vspace*{-8mm}\noindent \Large \sf Preprint}\par
  \vspace*{6mm}\noindent{\scriptsize \sf{\bfseries \MakeUppercase{#1}}}}

\fancyhead[L]{}
\fancyhead[R]{}

\begin{document}

\articletype{Letter}

\title{The Spatial Structure of the Gravitational Partition}

\author{Stephen Nelson}

\affil{Field Dynamics Research}

\ead{stephen@fielddynamics.org}

\begin{abstract}
The companion letter~\cite{Nelson2026b} derived the
exchange-symmetric partition $\mu(x) = x/(x{+}1)$ between the
two tetrahedral chambers of the closed frame, where
$x = g/a_0$ is the dimensionless acceleration.  That derivation
was algebraic: it determined the partition function but not its
spatial form.  This letter presents the spatial structure.
Evaluated radially about the field axis, the partition produces
a continuously graded field whose inner throat is
Christoffel-dominated ($\mu \to 1$), whose outer disk is
torsion-dominated ($\mu \to 0$), and whose crossover ring at
$x = 1$ is the MOND transition scale $a_0$.  The
24~torsion components organise into eight spiral channels
(the $2K = 8$ null directions of the closed frame), each
carrying three transverse rotation planes ($D = 3$).  The
chiral holonomy of these channels, computed in the
supplementary note, defines a rotation axis and forces the
field geometry to be planar.  The radial acceleration relation
$g_{\mathrm{bar}} = g_{\mathrm{obs}}^2/(g_{\mathrm{obs}} +
a_0)$ follows directly from the partition with no free
parameters.  The perceived separation of baryonic and dark
components is an artefact of projecting a 64-component
connection onto a 40-component subspace.
\end{abstract}

% ==================================================================

\section{The partition in space}

The companion letter~\cite{Nelson2026b} proved that the
exchange symmetry between TetA and TetB, combined with
rational bijectivity, uniquely determines the partition of
the gravitational field response between the two tetrahedral
chambers:
\begin{align}
  \text{TetA:} \quad & \mu(x) = \frac{x}{x + 1},
  \label{eq:mu} \\[4pt]
  \text{TetB:} \quad & 1 - \mu(x) = \frac{1}{x + 1},
  \label{eq:tetb}
\end{align}
where $x = g/a_0$ is the dimensionless acceleration.  The
derivation used no field equation and no fitting.  It required
only the parity relation between TetA and
TetB~\cite{Nelson2026a} together with the rational-bijection
classification on the projective line.  The present letter does
not re-derive this result.  It takes the function as given and
asks: what does this partition look like when it is evaluated
radially and revolved about a field axis?

The answer is Figure~\ref{fig1}.

\begin{figure}[t]
\centering
\includegraphics[width=\columnwidth]{fig1_tetrad_partition}
\caption{The tetrad partition rendered as a diagnostic
decomposition.  The two surfaces are not separate physical
fields but the complementary partition functions $\mu(x)$
and $1 - \mu(x)$ displayed simultaneously over the same
radial domain.  Blue: the TetA field fraction (the
Christoffel sector retained by general relativity).  Red:
the TetB field fraction (the torsion sector removed by the
zero-torsion condition).  At every radius,
blue~$+$~red~$= 1$.  The white ring marks the crossover at
$x = 1$, where both chambers contribute equally
($\mu = 1/2$).  Eight spiral arms alternate between TetA
(blue) and TetB (amber), corresponding to the $2K = 8$ null
directions of the closed frame.}
\label{fig1}
\end{figure}

The vertical axis in Figure~\ref{fig1} is the field fraction,
not a physical height.  The blue spike at centre is TetA
carrying nearly the entire field response.  The red plateau at
large radius is TetB carrying nearly all of it.  Between them,
the white ring marks the handover.  The textbook decomposition
of the full connection into Christoffel plus contorsion
corresponds, at the level of field magnitudes, to the
partition into $\mu$ and
$1 - \mu$~\cite{Nelson2026b}.  Figure~\ref{fig1} is that
decomposition rendered spatially.

The white ring sits at $x = 1$, where $\mu = 1/2$.  This is
the locus where both chambers contribute equally to the
field response.  Inside the ring, the Christoffel sector
dominates: general relativity is accurate and Keplerian
decline holds.  Outside the ring, the torsion sector
dominates: the field response that general relativity
discards becomes the majority.  This crossover is the MOND
transition scale $a_0$~\cite{Milgrom1983}.  It is not a
fitted parameter.  It is the geometric midpoint of the
exchange-symmetric partition, the point where the odds ratio
$\Omega = \mu/(1-\mu)$ equals unity.

The companion letter~\cite{Nelson2026b} presented the three
acceleration regimes in tabular form: strong field
($x \gg 1$, $\mu \approx 1$), transition ($x = 1$,
$\mu = 1/2$), and weak field ($x \ll 1$, $\mu \approx 0$).
Figure~\ref{fig1} gives those three regimes spatial identity:
the throat, the ring, and the disk.

The radial acceleration
relation~\cite{McGaugh2016,Li2018} is this partition,
rendered as geometry rather than as a scatter plot.  The blue
spike is the baryonic bulge.  The red disk is what standard
cosmological analysis attributes to an invisible mass
component.  They sum to one everywhere.

% ==================================================================

\section{The continuous field}

Figure~\ref{fig1} is explanatory: it separates the partition
into two complementary surfaces so that the reader can see
where each chamber dominates.  That separation is what a
torsion-free observer infers, one component they can account
for and one they cannot.  Figure~\ref{fig2} is ontological:
it shows the physical field as it is.  The gravitational field
is not two objects.  It is one.

\begin{figure}[t]
\centering
\includegraphics[width=\columnwidth]{fig2_tetrad_vortex}
\caption{The same partition rendered as one continuous
gravitational field.  A single surface whose colour encodes
the shift from TetA/Christoffel dominance (blue, inner
throat) to TetB/torsion dominance (red, outer disk).  The
MOND crossover ring (white) sits at $\mu = 1/2$.  The same
eight spiral channels from Figure~\ref{fig1} ride on one
surface, distinguished by colour alone.  There is no
boundary between regimes.  This is the physical geometry;
Figure~\ref{fig1} is its diagnostic decomposition.}
\label{fig2}
\end{figure}

The gravitational field is a single object with $K^3 = 64$
independent connection
components~\cite{Nelson2026b,Hehl1976}.  General relativity's
torsion-free condition projects this onto the 40~symmetric
components (Christoffel), discarding the 24~antisymmetric
components (torsion).  Figure~\ref{fig2} restores them.  The
result is one continuous surface with no separation, no
boundary, and no two-component model.

The shape is a steep inner throat opening into a flat extended
disk.  The throat is blue: high acceleration,
Christoffel-dominated.  The disk is red: low acceleration,
torsion-dominated.  The transition between them is smooth and
continuous.  This is the generic morphology of a spiral
galaxy~\cite{FamaeyMcGaugh2012}: a bright central bulge
surrounded by an extended disk.

The colour gradient in Figure~\ref{fig2} is not a rendering
choice.  It is governed by the logistic equation that the
partition satisfies~\cite{Nelson2026b}:
\begin{equation}
  x\,\frac{d\mu}{dx} = \mu(1 - \mu).
  \label{eq:logistic}
\end{equation}
Under the substitution $t = \ln x$, this becomes
$d\mu/dt = \mu(1-\mu)$: the canonical two-sector transition
law.  The product $\mu(1-\mu)$ on the right-hand side measures
the simultaneous presence of both tetrahedral chambers at a
given radius.  It is the continuous analogue of the octahedral
interaction core, the $s = 0$ state of the coupling
polynomial~\cite{Nelson2026b}.  The equation states that the
partition can shift only where both chambers contribute, which
is to say, only where the interaction core is active.  At the
deep centre, one chamber dominates and the partition is frozen.
At the rim, the same.  The transition peaks at the MOND ring
because that is where the core is fully realised:
$\mu(1 - \mu) = 1/4$.

The supplementary note to this letter establishes the
geometric origin of this constraint: every transition between
a TetA and a TetB null direction crosses a face of the
octahedral core, generating the chiral rotation that forces
the field into a disk.  The logistic equation is the
continuous, radial expression of that discrete crossing
requirement.

The perceived separation of baryonic and dark components in
Figure~\ref{fig1} is therefore an artefact.  It is the
projection of a single 64-component field onto a 40-component
subspace.  The observed gravitational discrepancy at galactic
scales is not in dispute: both the cold dark matter framework
and the present partition reproduce the radial acceleration
relation and the function $\mu(x) = x/(x{+}1)$.  They differ
on what the discrepancy represents.  The cold dark matter
framework interprets it as an invisible mass component, which,
having no geometric structure, distributes as a spherical
halo.  The closed tetrad frame identifies it as the field
response of the second tetrahedral chamber, carried by the
24~torsion components.  This response is not mass.  It is
topology.  Its spatial distribution is not fitted but
determined: the logistic equation~(\ref{eq:logistic}) governs
the radial profile, and the result is the continuously graded
disk of Figure~\ref{fig2}.

% ==================================================================

\section{The eight-channel structure}

Figures~\ref{fig1} and~\ref{fig2} both display eight spiral
arms, four coloured blue (TetA) and four amber (TetB),
alternating around the disk.  This section identifies what they
are and why they force the field into a disk rather than a
sphere.

The companion letter~\cite{Nelson2026b} proved (Section~5,
Theorem) that the 24~torsion components admit the
factorisation
\begin{equation}
  24 \;=\; \underbrace{8}_{\text{null directions}}
        \;\times\; \underbrace{3}_{\text{rotation planes}}
  \;=\; 2K \times D
  \label{eq:24}
\end{equation}
exclusively at $K = 4$.  In the companion letter's polynomial
language, $2K$ is the coupling range
$\Delta f = f(K,{+}1) - f(K,{-}1)$.  The algebraic identity
was established there.  The supplementary note to the present
letter identifies the two factors geometrically.  The
factorisation $8 \times 3$ is not one of several ways to
decompose 24; it is the unique decomposition in which one
factor is the vertex count of the closed frame and the other
is the spatial dimension.

The factor $2K = 8$ is the count of null directions of the
closed frame: the eight vertices of the stellated
octahedron~\cite{Nelson2026a}, equivalently the eight vertices
of the cube dual to the octahedral core.  Each face of the
octahedral core carries an attached tetrahedron whose apex is
one of these eight vertices.  The base face, shared with the
octahedron, is the propagation face of the corresponding
sub-vortex.  The three lateral faces are the three transverse
rotation planes.  The eight spiral arms in
Figures~\ref{fig1}--\ref{fig2} are these eight null directions
projected onto the disk.  Four belong to TetA.  Four belong to
TetB.  They alternate because adjacent faces of the octahedral
core always carry opposite chirality.

The supplementary note proves that parallel transport through
all eight null directions produces trivial holonomy when the
TetA/TetB chirality distinction is ignored, but nontrivial
holonomy when it is respected.  All twelve Hamiltonian cycles
on the cube graph were computed.  In the uniform case (all
rotations $+\theta$), all twelve yield the identity.  In the
chiral case (rotation sign determined by the source vertex's
sub-frame), all twelve yield nontrivial rotations: six at
$135^{\circ}$ and six at $155.8^{\circ}$.  The nontrivial
holonomy defines a rotation axis.  The field geometry is
planar, perpendicular to this axis.  Because the eight
channels all carry transport through the common field origin
with a consistent rotational sense, the holonomy does not
merely select an axis but organises the field into a
rotationally coherent disk orthogonal to it.  This is why the
partition forms a disk and not a sphere.  The axis is not imposed by
hand.  It emerges from the chirality of the closed frame.

The kinematic consequence follows directly from the partition.
At dimensionless acceleration $x = g_{\mathrm{obs}}/a_0$, the
baryonic acceleration is
$g_{\mathrm{bar}} = \mu\,g_{\mathrm{obs}}$, giving
\begin{equation}
  g_{\mathrm{bar}} =
  \frac{g_{\mathrm{obs}}^2}{g_{\mathrm{obs}} + a_0}.
  \label{eq:rar}
\end{equation}
At high acceleration ($x \gg 1$):
$g_{\mathrm{bar}} \approx g_{\mathrm{obs}}$, Keplerian
decline, general relativity is accurate.  At low acceleration
($x \ll 1$):
$g_{\mathrm{bar}} \approx g_{\mathrm{obs}}^2/a_0$, flat
rotation curves.  The radial acceleration relation established
empirically by McGaugh, Lelli \&
Schombert~\cite{McGaugh2016}, with 0.057~dex scatter across
175~individual galaxies~\cite{Li2018} and no credible
variation in $a_0$, is reproduced by Eq.~(\ref{eq:rar}) with
no free parameters.  The full 179-galaxy rotation curve
analysis is presented in a companion
paper~\cite{Nelson2026c}.

% ==================================================================

\section{What the partition resolves}

The spatial partition developed in Sections~1--3 makes three
observational predictions that distinguish it from the cold
dark matter framework.

\textit{Disk geometry.}  The cold dark matter framework
distributes the invisible component as a spherical halo
because structureless mass has no preferred
direction~\cite{NFW1996}.  The closed tetrad frame distributes
it as a disk because the chiral holonomy of the eight-channel
structure defines an axis (supplementary note).  Observed
spiral galaxies are disks.  This is a zero-parameter geometric
prediction.

\textit{Bounded central density.}  The TetB field fraction
$1/(x{+}1)$ is bounded above by unity for all $x$.  The
torsion-sector field acceleration never exceeds $a_0$.  The NFW
profile diverges as $r \to 0$~\cite{NFW1996}.  The partition's
boundedness resolves the cusp-core tension without baryonic
feedback or self-interacting dark matter.  The resolution is
algebraic: $1/(x{+}1) < 1$ is a property of the constitutive
law, not of any galaxy model.

\textit{No new particles.}  The gravitational discrepancy is
not mass.  It is the field response of the second tetrahedral
chamber, carried by the 24~torsion components that the
zero-torsion condition removes from the connection.  The
missing component is not in the matter sector.  It is in the
connection.

An alternative relativistic approach to the MOND phenomenology
has been developed by Skordis \&
Z\l{}o\'{s}nik~\cite{SkordisZlosnik2021}, who introduce
additional vector and scalar fields to obtain a theory
compatible with the cosmic microwave background power spectrum.
The present framework adds no fields.  It restores the
24~connection components that general relativity's zero-torsion
condition removes, recovering the complete geometry of the
tetrad frame.

The torsion-free specialisation adopted in 1915 discards the
24~antisymmetric components of the affine connection.  The
field response they carry is precisely what, in their absence,
must be attributed to invisible mass.

% ==================================================================

\section{Data and reproducibility}

The chiral holonomy computation is provided as a supplementary
note (\texttt{supplementary\_chiral\_holonomy.pdf}).  A
self-contained verification script
(\texttt{holonomy\_verify.py}, approximately forty lines of
Python using NumPy) reproduces the holonomy of all twelve
Hamiltonian cycles on the cube graph for both uniform and
chiral transport.  Interactive renderings of
Figures~\ref{fig1} and~\ref{fig2} are provided as
supplementary HTML files permitting rotation, zoom, and
background toggling.

% ==================================================================

\begin{thebibliography}{99}

\bibitem{Nelson2026a}
S.~Nelson,
\textit{The Coordinate-Space Geometry of the Tetrad Frame},
Class.\ Quantum Grav.\ (2026), submitted (CQG-115428).

\bibitem{Nelson2026b}
S.~Nelson,
\textit{The Other Half of Gravity: Torsion and the
Two-Chamber Structure of the Tetrad Frame},
Class.\ Quantum Grav.\ (2026), submitted (CQG-115447).

\bibitem{Nelson2026c}
S.~Nelson,
\textit{Mass Decoding from the Scalar-Tensor Transition},
Mon.\ Not.\ R.\ Astron.\ Soc.\ (2026), submitted.

\bibitem{McGaugh2016}
S.~S.~McGaugh, F.~Lelli and J.~M.~Schombert,
\textit{Radial Acceleration Relation in Rotationally
Supported Galaxies},
Phys.\ Rev.\ Lett.\ \textbf{117} (2016) 201101.

\bibitem{Li2018}
P.~Li, F.~Lelli, S.~McGaugh and J.~Schombert,
\textit{Fitting the radial acceleration relation to
individual SPARC galaxies},
Astron.\ Astrophys.\ \textbf{615} (2018) A3.

\bibitem{Milgrom1983}
M.~Milgrom,
\textit{A modification of the Newtonian dynamics as a
possible alternative to the hidden mass hypothesis},
Astrophys.\ J.\ \textbf{270} (1983) 365.

\bibitem{FamaeyMcGaugh2012}
B.~Famaey and S.~S.~McGaugh,
\textit{Modified Newtonian Dynamics (MOND): Observational
Phenomenology and Relativistic Extensions},
Living Rev.\ Relativity \textbf{15} (2012) 10.

\bibitem{Hehl1976}
F.~W.~Hehl, P.~von~der~Heyde, G.~D.~Kerlick and
J.~M.~Nester,
\textit{General relativity with spin and torsion:
Foundations and prospects},
Rev.\ Mod.\ Phys.\ \textbf{48} (1976) 393.

\bibitem{NFW1996}
J.~F.~Navarro, C.~S.~Frenk and S.~D.~M.~White,
\textit{The structure of cold dark matter halos},
Astrophys.\ J.\ \textbf{462} (1996) 563.

\bibitem{SkordisZlosnik2021}
C.~Skordis and T.~Z\l{}o\'{s}nik,
\textit{New Relativistic Theory for Modified Newtonian
Dynamics},
Phys.\ Rev.\ Lett.\ \textbf{127} (2021) 161302.

\end{thebibliography}

\end{document}
